% Folder 74, batch 02 (archivists' pages 21-40).
% Pass: Opus 5.5 (claude-opus-5-5), 2026-10-03 — first pass, unchecked against the pages by a human.
%
\documentclass[11pt,a4paper]{article}
\input{../preamble/grothendieck}

\folder{74}
\batch{2}
\pages{21}{40}
\foldertitle{Complexes cubiques : notes manuscrites (s.d.).}
\dating{[à partir de 1976]}
\watermark{Édition de démonstration}

\begin{document}

\page{21}
\note{feuillet de brouillon couvert de notes éparses, sans phrase suivie ; on en donne les formules et les mots lisibles, dans l'ordre de la page, de haut en bas.}

$B \quad B' \amalg B''$

\[
\begin{matrix}
\xi_{14} & \xi_{15} & \xi_{16}\\
\xi_{24} & \xi_{25} & \xi_{26}\\
\xi_{34} & \xi_{35} & \xi_{36}
\end{matrix}
\]

\note{à droite, une multiplication posée : $36 \times 36 = 1296$, puis le nombre $1440$ inscrit sous le résultat.}

… triades $\longleftrightarrow$ hexagones

\[
6^4 = 36 \times 36 = 2^4 3^4, \qquad 720 = 2^4 \cdot 3^2 \cdot 5, \qquad 9/5
\]
\note{l'exposant $4$ de $2$ est récrit sur un $3$.}

\note{dans la marge droite, une accolade réunit les nombres $40$, $36$ (récrit sur un $45$ biffé), $45$.}

\uncertain{Partition} \uncertain{ordonnée} d'une \ill{}

\[
\xi_{12}\ \xi_{23}\ \xi_{31} \;\Big|\; \xi_{45}\ \xi_{56}\ \xi_{64}
\]

\drawing{un prisme hexagonal (ou un cube tronqué) en perspective ; au-dessous, un graphe en éventail : un sommet relié par des arcs à cinq sommets alignés, et un petit trièdre de trois segments issus d'un point.}

\note{dans la moitié inférieure, d'autres fragments, dont une partie écrite tête-bêche : $\xi_{ij}$ (deux fois), $\{\xi_1, \ldots, \xi_6\}$, $\xi_{12}\ \xi_{13} - \xi_{16}$, et deux lignes de mots illisibles.}

\page{22}
\drawing{deux versions d'un même schéma. En haut, trois blocs de symboles reliés par des faisceaux de traits croisés. Le premier bloc, surmonté de « $\{1,2,3\} = I_3$ », porte les trois lignes $\xi_1\ \xi_2\ \xi_3$ / $\hat\xi_1\ \hat\xi_2\ \hat\xi_3$ / $\xi_{23}\ \xi_{31}\ \xi_{12}$ ; le deuxième, surmonté de « $\{4,5,6\} =$ » et d'un symbole encadré $\tau_3$, porte $\xi_4\ \xi_5\ \xi_6$ / $\hat\xi_4\ \hat\xi_5\ \hat\xi_6$ / $\xi_{56}\ \xi_{64}\ \xi_{45}$, avec en regard l'accolade « $I_1 = \{$\ill{}$, *, !\}$ » ; le troisième, en bas, porte $\xi_{14}\ \xi_{15}\ \xi_{16}$ / $\xi_{24}\ \xi_{25}\ \xi_{26}$ / $\xi_{34}\ \xi_{35}\ \xi_{36}$, avec « $\{1,2,3\} = I_3$ » à droite, « $\{4,5,6\}$ » au-dessous et un \uncertain{$\tau_3$} cerclé. À droite, un triangle dont les trois sommets portent chacun un des trois blocs, marqué « I ». En bas de page, le même schéma redessiné : le bloc $\xi_1 \ldots \xi_{12}$ relié d'un côté au bloc $\xi_{4} \ldots \xi_{45}$, de l'autre au bloc $\xi_{14} \ldots \xi_{36}$.}

\page{23}
\marginal{$q$, $q'$, $a$ fixes ; $p$, $p'$ échangés}

$p\,q\,p'\,q'$ \qquad
\note{à côté, un carré pointillé aux sommets $p$, $q$ (en haut), $q'$, $p'$ (en bas).}

\drawing{un polyèdre en perspective : en haut le sommet $q$, puis $q'$ ; au milieu un quadrilatère $s$, $r'$, $s'$, $r$, dont les côtés sont fléchés ($s \to r' \to s' \to r \to s$) ; au centre le point $a$ ; en bas les sommets $p$, $p'$ ; les arêtes joignent $q$, $q'$ au quadrilatère, et le quadrilatère à $p$, $p'$.}

\struck{\ill{}}

$x$ et $y$ (distincts de $a$, $q$ \uncertain{et $q'$} \ill{}) sont liés exactement \uncertain{par deux} \ill{}, l'une \uncertain{à} $(r, a, r')$, l'autre : $(s\, a\, s')$.

Montrons que \add{si} $x, y$ sont liés, ils \ill{} \struck{\ill{} liés, i.e.} \ill{} $=$ \uncertain{relatifs}.

$x, y \in T'$, $x \neq y$ : pour qu'ils soient liés, il faut et il suffit qu'ils aient des arêtes distinctes, et \uncertain{convergentes} \ill{} à des sommets distincts sur $T$.

\drawing{au bas de la page, quatre figures : un prisme aux sommets $c$, $b$, $t$, $t'$, $s$, $s'$, $a$, $a'$ ; un triangle dans lequel sont emboîtés deux angles de même sommet ; un double cône de sommets $s$ (en haut) et $s'$ (en bas), dont la ceinture porte $a$, $b'$, $c$, $a'$, $c'$, $b$, à côté d'un carré de $3 \times 3$ points ; un autre double cône analogue, de sommets $s'$ et $s$, portant $a$, $b'$, $c'$, $b$, $c$, $a'$ et des points $\alpha$, $\beta$, $\gamma$ à l'intérieur. Une croix au bas de la page.}

\[
\begin{matrix}
(\alpha b, \alpha c) & (\alpha b', \alpha c')\\
(\beta c, \beta a) & (\beta c', \beta a')\\
(\gamma a, \gamma b) & (\gamma a', \gamma b')
\end{matrix}
\]

\page{25}
\note{page de dénombrements, en colonnes séparées par des traits ; on suit l'ordre de la colonne de gauche, puis celle de droite.}

\[
\Gamma_0 = 1 \begin{pmatrix} 27 \\ 36 \\ 40 \end{pmatrix}
\]
\note{le $\Gamma_0$ pourrait être le $r_0$ des pages suivantes. Les trois nombres sont entourés ; à droite de la parenthèse, les chiffres \uncertain{$9$ et $4$}.}

\struck{$720 \cdot 6$} \qquad $9 \cdot 8 \cdot 6$ \qquad $72 \cdot 37$ \qquad $\dfrac{72 \cdot 71}{2} =$

\[
72 \cdot 36\tfrac{1}{2} = 72 \cdot \tfrac{71}{2} + 72
\]

\uncertain{3-baies} \ill{} \qquad
\uncertain{paires de \uncertain{voisins} à $120^\circ$} \quad $\Big|\ 720,\ \uncertain{4}\,720$

$720$ \uncertain{groupés} en $120$ hexagones qui se groupent par triples en $40$ trihexagones de triades

\[
3!\, 3!\, 2! = 72
\]

\uncertain{bitriade} $+$ partie à trois élt de l'ens \add{$B/$}$B_i$ \ill{}

\uncertain{$\Updownarrow$} partition \ill{} : trois él. $=$ triplets

$\Updownarrow$ \struck{\ill{}} \add{paires de deux} \uncertain{3-baies} d'\uncertain{intersection} ordonnés : $3$ él.

$\Updownarrow$ paire de voisins à $60^\circ$

\[
\{\xi_{ij}\ \xi_{jk}\ \xi_{ki}\ \hat\xi_i \ldots\}, \qquad \{\xi_{ij}\ \xi_{jk}\ \xi_{ki}\ \hat\xi_i\ \hat\xi_j \ldots\}
\]

\marginal{$(\xi_i, \xi_i^{\wedge})$ \quad $\{i, j, k\}$ \quad $\{\xi_{ij}, \xi_{jk}, \xi_{ki}\}$ ; \uncertain{stabilisateur} $\mathfrak{S}_3 \times \mathfrak{S}_3 \times \mathfrak{S}_2$ d'ordre \struck{$72$} $= 2^3 3^2$ ; il y en a $2^4 \cdot 3^2 \cdot 5 = 720$.}
\note{la ligne « Ens$_3 \times$ Ens$_3 \times$ Ens$_2$ » est écrite sous ce calcul.}

\uncertain{feuille} 3-couverte d'él.\ de $\Delta$

$\big|$ hexagone distingué de triades

$\big|$ syst.\ de trois bitriades \struck{\ill{} position \ill{}} \add{\uncertain{en} position \ill{}}

$\big|$ plans de deux voisins \uncertain{non} orthogonales

\[
\begin{matrix}
\xi_{ip} & \xi_{jp} & \xi_{kp}\\
\xi_{iq} & \xi_{jq} & \xi_{kq}\\
\xi_{ir} & \xi_{jr} & \xi_{kr}
\end{matrix}
\]

\marginal{\uncertain{stabilisateur} $(\mathfrak{S}_3 \times \mathfrak{S}_3 \times \mathfrak{S}_3) \cdot \mathfrak{S}_2$ \add{ext} d'ordre $6^3 \cdot 2 = 432$ ; il y en a $120$.}

\struck{\uncertain{décomposition de $\Delta$ en $3$ \uncertain{grilles} 3-couvertes}}

trihexagone distingué de triades \marginal{\uncertain{les} syst.\ distingués à $18$ triades \add{\uncertain{$9$ couples} de triades \ill{}}}

décomposition \add{orth.} (distinguée) de $E$ en somme \add{orth.} de trois plans de voisins

syst.\ \struck{de} \add{distingués} $9$ bitriades (en $9$ \uncertain{paires} de voisins opposés)

\marginal{\uncertain{stabilisateur} $\mathfrak{S}_3 \cdot (\mathfrak{S}_3 \times \mathfrak{S}_3 \times \mathfrak{S}_3)$ ; \uncertain{Cel.}\ des \uncertain{ens} : $9$ éléments munis d'une partition en trois parties équipotentes.}

\[
\frac{6 \cdot 5 \cdot 4}{3!\, 2}
\]

\page{26}
\note{page de brouillon où se mêlent tableaux de symboles, dénombrements et petits graphes ; un long trait oblique la traverse sans la biffer. On transcrit les formules dans l'ordre de lecture et l'on décrit les figures.}

\[
\begin{matrix}
\xi_1 & \xi_2 & \xi_3 & & \xi_4' & \xi_5' & \xi_6'\\
\xi_1' & \xi_2^{*} & \xi_3' & & \xi_4^{*} & \xi_5^{*} & \xi_6^{*}\\
\xi_{23}^{*} & \xi_{31}^{*} & \xi_{12}^{*} & & \xi_{56} & \xi_{64} & \xi_{45}
\end{matrix}
\]
\note{les indices supérieurs (accents, étoiles) sont lus avec doute ; entre les deux blocs, un faisceau de traits croisés, en partie raturé, relie les lignes.}

\[
\begin{matrix}
\xi_{14} & \xi_{15} & \xi_{16}\\
\xi_{24} & \xi_{25} & \xi_{26}\\
\xi_{34} & \xi_{35} & \xi_{36}
\end{matrix}
\]
\note{ce troisième bloc est relié aux deux premiers par des arcs ; à gauche « \struck{$P$} $\times\, Q$ », avec au-dessous $\{b, b', b''\}$ et $\{1, 2, 3\}$ ; à droite « $P \times R$ », avec au-dessous \uncertain{$\{b, b', b''\}$} et $(4, 5, 6)$ ; sous le bloc, « $R \times P$ ».}

\marginal{Colonne de droite, de haut en bas :
$1, 10, 16$ ;
$2, \struck{5}, 10, 10$ ;
$1, 1, 5, 5, 5, 10$ ;
$3, 3, 6, 6, 9$ ;
\struck{\ill{}} $1, 1, 1, 2, 2, 2, 3, 3, 3, 3, 6$ ;
$1, 2, 2, 4, 4, 6, 8$ ;
$15$ fois $1$, $6$ fois $2$, \struck{\ill{}} ;
$1, 1, 5, 5, 5, 10$ (relié par un trait à $(4,5,6)$) ;
$27$ fois $1$ ;
\uncertain{$1, 6, 25$} ;
$27$ fois $1$ ;
\uncertain{comme $3^{\circ}$)} ;
$9, 9, 9$ ;
$27$.}

$\{1,2,3\} \times \{b, b', b''\}$, \quad $\{4,5,6\} \times \{b, b', b''\}$, \quad $\{1,2,3\} \times \{4,5,6\}$

$(E_\lambda)_{\lambda \in \Lambda}$, \quad card $E_\lambda = 3$, \quad card $\Lambda = 3$

$E, E', E''$

\[
E \qquad \mathfrak{S}_3 \times \mathfrak{S}_3 \times \mathfrak{S}_3 \subset W
\]

\[
2^3 3^3, \qquad 2^4 \cdot 3 \cdot 5 = 240
\]

\[
r_6 = \eta - \xi_1 - \xi_2 - \xi_3, \qquad r_0 = 2\eta - \xi
\]
\note{lecture de $\eta$ et $\xi$ incertaine ; la lettre $r$ (racine), qui revient p.~27, ressemble ici à un $\Gamma$.}

\[
E \times E' \amalg E' \times E'' \amalg E'' \times E
\]

$(x, x')$ et $(y', y'')$ liés ssi $x' \neq y'$

$(x, x')$ et $(y, y')$ liés ssi $x \neq y$ \uncertain{et} $x' \neq y'$

\drawing{une petite figure en angle, puis deux étoiles de segments reliées par une arête ; au-dessous un faisceau de traits, une « maison » (carré surmonté d'un toit, un triangle inscrit, un axe pointillé), un segment divisé en trois, un graphe biparti de trois sommets en haut et trois en bas reliés par de nombreuses arêtes, et plusieurs petits schémas de trois traits horizontaux reliés par des obliques, l'un d'eux étiqueté $\alpha$, $\beta$, $\alpha$.}

\[
(2\eta - \xi) \cdot \xi_j = 2
\]
\note{à côté : « $\xi_{1}$ », « $1\ 1$ », et une suite illisible entre accolades.}

\marginal{écrit verticalement dans la marge gauche : $27, 36, 40$ ; $\xi_j = 3\eta - \xi \ldots$ ; $2\eta - \xi$, $2\eta - \xi_i - \ldots$}

$\mathfrak{S}_3 \quad \sim \qquad 6^4 = 2^4 3^4$ \quad $40$ (entouré)

\[
F \to E, \qquad \coprod_{i \in E}\ \prod_{j \in E - \{i\}} F_j
\]

\[
81 \times 16 = 1296
\]

\ill{} on \uncertain{particulier} \struck{\ill{}} $K$ \struck{telle que} de card.~$2$ \ill{} qui se \ill{} dans une \uncertain{fibre}

\[
9 \times 9 = 81, \qquad 3 \times 3 \times 3 = 27, \qquad 54
\]

\[
\frac{27 \cdot 16 \cdot 10}{6} = 9 \cdot 16 \cdot 5 = 720
\]
\note{écrit verticalement le long du bord droit ; le dénominateur et le $16$ du second membre sont lus avec doute.}

\page{27}
\note{trois figures de même forme, l'une sous l'autre : six groupes de trois symboles disposés aux sommets d'un hexagone (deux triangles en pointillé formant une étoile), deux sommets voisins hachurés ; chaque côté porte une racine. On transcrit pour chaque figure les sommets, dans l'ordre haut, haut droite, bas droite, bas, bas gauche, haut gauche, puis les étiquettes.}

\textbf{Première figure.} Sommets : $\xi_1\,\xi_2\,\xi_3$ ; $\xi_4\,\xi_5\,\xi_6$ ; $\xi_{23}\,\xi_{31}\,\xi_{12}$ ; $\xi_4'\,\xi_5'\,\xi_6'$ ; $\xi_1'\,\xi_2'\,\xi_3'$ ; $\xi_{56}\,\xi_{64}\,\xi_{45}$. Côtés : $r_0$, $-r_0$, $r_6$, $-r_{123} = -r_6$, $-r_{456}$, avec
\[
r_{456} = 3(r_1 + r_2 + r_3) + 2r_4 + r_5 + r_6
\]
\note{le coefficient de $r_4$ et le signe devant $r_3$ sont lus avec doute.}

\textbf{Deuxième figure.} Sommets : \uncertain{$\xi_{1}\,\xi_{1}'\,\xi_{13}$} ; $\xi_{24}\,\xi_{25}\,\xi_{26}$ ; $\xi_3\,\xi_3'\,\xi_{12}$ ; $\xi_{14}\,\xi_{15}\,\xi_{16}$ ; $\xi_2\,\xi_2'\,\xi_{31}$ ; $\xi_{34}\,\xi_{35}\,\xi_{36}$. Côtés :
\[
r_{2,1} = -r_{12} = -r_1, \quad r_{2,3} = r_2, \quad r_{1,3} = r_1 + r_2, \quad r_{1,2} = r_1, \quad r_{3,2} = -r_2, \quad r_{3,1} = -r_1 - r_2 .
\]

\textbf{Troisième figure.} Sommets : $\xi_4\,\xi_4'\,\xi_{56}$ ; $\xi_{15}\,\xi_{25}\,\xi_{35}$ ; $\xi_6\,\xi_6'\,\xi_{45}$ ; $\xi_{14}\,\xi_{24}\,\xi_{34}$ ; $\xi_5\,\xi_5'\,\xi_{64}$ ; $\xi_{16}\,\xi_{26}\,\xi_{36}$. Côtés :
\[
r_{5,4} = -r_{45} = -r_4, \quad r_{5,6} = r_5, \quad r_{4,6} = r_4 + r_5, \quad r_{4,5} = r_4, \quad r_{6,5} = -r_5, \quad r_{6,4} = -r_4 - r_5 .
\]

\marginal{En regard des trois figures, dans la marge droite : $[r_6, r_0] = [r_6, r_{456}] = [r_{123}, r_{456}]$ ; $[r_1, r_2]$ ; $[r_4, r_5]$.}

$\overbrace{r_1\ r_2}\ \overbrace{r_4, r_5}\ \ r_0, r_6$ \ engendrent un sous-groupe d'indice $3$ dans \uncertain{$W$}

\struck{$[r_1, r_2]$ $[$}

\begin{tikzcd}[column sep=small, row sep=small]
r_1 \arrow[r, no head] & r_2 \arrow[r, no head] & r_3 \arrow[r, no head] \arrow[d, no head] & r_4 \arrow[r, no head] & r_5\\
 & & r_6 & &
\end{tikzcd}

\note{le diagramme de Dynkin de type $E_6$, tel qu'il est dessiné : $r_1 - r_2 - r_3 - r_4 - r_5$, avec $r_6$ attaché à $r_3$.}

\[
\begin{matrix}
\xi_1 & \xi_2 & \xi_3 & \text{---} & \xi_4' & \xi_5' & \xi_6'\\
\xi_1' & \xi_2' & \xi_3' & \text{---} & \xi_4 & \xi_5 & \xi_6\\
\xi_{23} & \xi_{31} & \xi_{12} & \text{---} & \xi_{56} & \xi_{64} & \xi_{45}
\end{matrix}
\]
\note{au-dessous, le bloc $\xi_{14}\ \xi_{15}\ \xi_{16}$ / $\xi_{24}\ \xi_{25}\ \xi_{26}$ / $\xi_{34}\ \xi_{35}\ \xi_{36}$, relié par des arcs aux deux blocs précédents, comme p.~26.}

\page{28}
\note{un long trait courbe, au crayon, traverse la page de haut en bas sans biffer le texte.}

\textbf{Prop.} Les \struck{bitriades} \add{\uncertain{vrais} 3-bisons} ordonnées sont conjuguées. Si $I$ est une telle partie, $\exists$ exactement deux bases qui la contiennent, soit $I$ et $I'$, et $(\underbrace{I - J}_{K}) \cup (\underbrace{I' - J}_{K'})$ est de cardinal $6$, et formé de l'ens.\ des \uncertain{$\delta \in \Delta - J$} qui ne sont liés à aucun él.\ de $K$. De plus, $\exists !$ \uncertain{bitriade} $I \cup I'$ contenant $K \cup K'$, et $K$, $K'$ sont \uncertain{respectivement} les \uncertain{intersections} de $I$, $I'$ avec $K \cup K'$, \add{$J$ donc $K \cup K'$ \ill{}} et \struck{\ill{}} \struck{correspondance} par $\sigma$ (l'unique \uncertain{automorphisme} $\neq \mathrm{id}$ qui \add{\uncertain{respecte}} \struck{\ill{}} \uncertain{fixe} la \uncertain{bitriade} \ldots) et \uncertain{si}

\[
J = \{\lambda_{45}, \lambda_{56}, \lambda_{61}\}, \qquad
K \cup K' = \{\underbrace{\xi_1, \xi_2, \xi_3}_{K}, \underbrace{\xi_1', \xi_2', \xi_3'}_{K'}\}
\]

\[
6! = 720 = 2^4 3^2 5, \qquad 6^4 = 2^4 3^4 = 1296
\]

$72$ bases \ $\big/$ \ $36$ \uncertain{bitriades} \quad $40$ trihexagones

\note{à droite, la multiplication posée $36 \times 36 = 1296$ (avec les produits partiels $216$ et $108$).}

\[
\begin{array}{ll}
\text{sommets} & 27\\
\text{\uncertain{bitriades}} & 36\\
\text{trihexagones} & 40\\
\text{triangles} & 45
\end{array}
\]

\drawing{un triangle épais d'où rayonnent de nombreux triangles plus fins, en éventail, suivi d'une accolade.}

$\mathfrak{S}_4$ \quad $\mathbb{F}_2^4$ \quad card.\ $24 \cdot 16$

\uncertain{quadriques} affines de dim.~$3$ sur $\mathbb{F}_2$ \qquad $8 \cdot \mathrm{Gl}(3, \mathbb{F}_2)$

\[
\text{\struck{$2$}} \quad \mathfrak{S}_?\,/\,\mathfrak{S}_4 \qquad \mathbb{F}_2^{3}
\]
\note{la ligne est surchargée : un $\mathfrak{S}_4$ est écrit au-dessus, l'indice du premier $\mathfrak{S}$ est illisible, et l'exposant de $\mathbb{F}_2$ a été corrigé ; le $7$ de la ligne suivante est récrit sur un autre chiffre.}

\[
72 \cdot 2^{4} \qquad 2^{7} \cdot 9 \qquad (\sim 16\ \text{él.})
\]

$A(i)$ \uncertain{ensemble} à $4$ él.\ \quad $i \in T$

$P(i)$ torseur sous $\mathbb{F}_2^{A(i)}$ \quad \uncertain{$k(i)$}

$\forall i$, bijection entre $P(i)$ et $\displaystyle\coprod_{j \in T - \{i\}} \ldots\ (\alpha \in A_j)$

\[
\frac{8 - 2}{2} = 3 \qquad \text{\struck{$P(i)$}} \quad \text{\struck{$2$}}
\]

\[
P(i) / \text{\struck{$V$}}\,V'(i) \simeq T - \{i\}
\]

\[
0 \to k \to V'(i) \to V(i) \to 0
\]
\note{sous la suite, les dimensions $1$, $2$, $1$.}

\[
(1 + q(1 + q^{-1}))\, q^3 (q-1)^2 \qquad q = 2
\]

\[
2^3 \cdot 3 \cdot 5 = 120 \ !!! \qquad \text{\struck{$P(i)$}}
\]

\drawing{une petite figure : un segment horizontal et une flèche verticale montant vers un point.}

$20$

$V$ \uncertain{vectoriel} de dim.~$3$ sur $\mathbb{F}_2$ \quad \struck{$V'(i)$} \quad $k \subset V$, \quad $\downarrow$ \quad $\mathrm{Aut}(E)$

\[
V \simeq k(\Delta_i)' \ \text{pour tout } i, \qquad \check{k} = k
\]

$A_3 \times$

\marginal{Dans la colonne de droite : $V'(i)$ ; $P_i$ ; $A_1$ ; $A_3$, $A_2$, \uncertain{3 ensembles à 4 él.} ; $V \supset k$ \uncertain{vectoriel} de dim.~$3$ sur $\mathbb{F}_2$.}

\page{29}
$\xi_1\ \xi_2\ \xi_3$ \quad $\xi_{32}$ \quad / \quad $\xi_1'\ \xi_2'\ \xi_3'$
\drawing{un segment divisé en deux.}

\[
\frac{72 \cdot 6!}{72 \cdot 6} = \frac{72 \cdot 720}{72 \cdot 6} = 120
\]

\[
\begin{matrix}
\xi_1\ \xi_2\ \xi_3 & \xi_4\ \xi_5\ \xi_6 & \xi_{12}\ \xi_{23}\ \xi_{34}\\
\xi_1'\ \xi_2'\ \xi_3' & \xi_4'\ \xi_5'\ \xi_6' & \xi_{45}\ \xi_{56}\ \xi_{61}
\end{matrix}
\]
\note{le dernier bloc est souligné ; les indices $34$ et $61$ (au lieu de $31$ et $64$) sont ceux de la page.}

\textbf{Hexagone.} Six groupes de trois symboles aux sommets d'un hexagone : en haut $\xi_{12}\ \xi_{23}\ \xi_{34}$ ; puis, dans le sens des aiguilles d'une montre, $\xi_4\ \xi_5\ \xi_6$, $\xi_1\ \xi_2\ \xi_3$, $\xi_{45}\ \xi_{56}\ \xi_{61}$ (en bas), $\xi_1'\ \xi_2'\ \xi_3'$, $\xi_4'\ \xi_5'\ \xi_6'$. Les côtés portent, dans le même ordre, $r_{456}$, $r_0$, $r_6$, $-r_{456}$, $-r_{\ill{}}$, $-r_6$.

\marginal{Dans un cadre, à droite :
triades $\longleftrightarrow$ hexagones ;
hex.\ triades $\longleftrightarrow$ bitriangles ;
$\updownarrow$ \quad $\updownarrow$ ;
plan $A_2$ $\longleftrightarrow$ bloc de triades ;
trihexatriades $\longleftrightarrow$ \struck{\ill{}} \add{triples de} bitriangles ;
$\updownarrow$ \quad $\updownarrow$ ;
décomposition \struck{$A_2 \ldots A_2$} \add{triple $A_2$} $\longleftrightarrow$ triple \struck{\ill{}} bloc de triades.
Au-dessous : $\xi_{ij}$, $1 \leq i \leq 3$, $4 \leq j \leq 6$.}

\[
\begin{matrix}
\xi_{14} & \xi_{15} & \xi_{16}\\
\xi_{24} & \xi_{25} & \xi_{26}\\
\xi_{34} & \xi_{35} & \xi_{36}
\end{matrix}
\]

NB le \uncertain{nombre} de $3$ bases (\uncertain{ou} couples de voisins \uncertain{à} $60^\circ$)
\[
36 \cdot \binom{6}{3} = 720, \qquad \binom{6}{3} = \frac{6 \cdot 5 \cdot 4}{3!} = 20
\]
\note{à droite de $\binom{6}{3}$, une petite fraction surchargée ($\tfrac{4}{24}$ ?), et le $720$ est récrit sur un autre nombre.}

nb de paires de \uncertain{voisins} faisant un \uncertain{angle} de $60^\circ$

\[
\text{plan } A_2 : \quad \frac{720}{6} = 120
\]

\struck{plan $A_2$}

sous-syst.\ $A_2$ $\updownarrow$ $3$ blocs de \uncertain{bitriades} $\updownarrow$ \uncertain{bitriangle} $\updownarrow$ hexagone de triades
\note{ces quatre termes sont écrits en colonne, reliés par des doubles flèches verticales, et réunis par une accolade qui renvoie aux deux groupes ci-dessous.}

\[
\mathfrak{S}_2 \cdot (\mathfrak{S}_3 \times \mathfrak{S}_3), \qquad
\mathfrak{S}_3 \cdot (\mathfrak{S}_3 \times \mathfrak{S}_3 \times \mathfrak{S}_3)
\]

\[
\begin{matrix}
\xi_{14} & \xi_{25} & \xi_{36}\\
\xi_{26} & \xi_{34} & \xi_{15}\\
\xi_{35} & \xi_{16} & \xi_{24}
\end{matrix}
\]

\[
\left\{
\begin{array}{ll}
\text{sommets} & 27\\
\text{triangles} & 45\\
\text{bibases} & 36\\
\text{trois-neufs} & 40
\end{array}
\right.
\]
\drawing{à gauche du tableau, trois petits symboles : un triangle surmonté d'un point, une ligne de points, et un carré quadrillé marqué « \uncertain{triple} ».}

\page{30}
\note{trois lignes, chacune ouverte par un crochet, de trois paires de triplets ; des traits relient certains symboles entre les lignes, étiquetés $\xi_1$, $\xi_2$, $\xi_3$, $\xi_4$, $\xi_5$, $\xi_6$. Le reste de la page est blanc.}

\[
\begin{array}{lll}
(\xi_1\,\xi_2\,\xi_3)(\xi_4'\,\xi_5'\,\xi_6') & (\xi_4\,\xi_5\,\xi_6)(\xi_1'\,\xi_2'\,\xi_3') & (\xi_{23}\,\xi_{31}\,\xi_{12})(\xi_{56}\,\xi_{64}\,\xi_{45})\\[4pt]
(\xi_1\,\xi_1'\,\xi_{23})(\xi_{14}\,\xi_{15}\,\xi_{16}) & (\xi_{24}\,\xi_{25}\,\xi_{26})(\xi_2\,\xi_2'\,\xi_{31}) & (\xi_3\,\xi_3'\,\xi_{12})(\xi_{34}\,\xi_{35}\,\xi_{36})\\[4pt]
(\xi_4\,\xi_4'\,\xi_{56})(\xi_{14}\,\xi_{24}\,\xi_{34}) & (\xi_5\,\xi_5'\,\xi_{64})(\xi_{15}\,\xi_{25}\,\xi_{35}) & (\xi_6\,\xi_6'\,\xi_{45})(\xi_{16}\,\xi_{26}\,\xi_{36})
\end{array}
\]
\note{plusieurs indices sont surchargés (en particulier dans $(\xi_{56}\,\xi_{64}\,\xi_{45})$ et $(\xi_1\,\xi_1'\,\xi_{23})$) ; lecture par endroits incertaine.}

\page{32}
$V$ groupe comm.\ (écrit additivement), $I$ ens.\ de cardinal $3$,
\[
\Gamma = \mathrm{Ker}(V^I \to V),
\]
$\omega$ ens.\ des \add{\uncertain{deux}} ordres circulaires sur $I$ (c'est un \struck{$\mathbb{Z}$} \uncertain{$(\mathbb{Z}/2)$}-tors.),
\[
V^\omega = V \wedge \omega
\]
(\struck{$\mathbb{Z}$} \uncertain{$\mathbb{Z}/2$} qui opère sur $V$ par chgt de signe).

Pour $i \in I$, \struck{\ill{}} \ill{}
\[
V^\omega \xrightarrow{\ \varphi_i\ } \Gamma, \qquad x \wedge \rho \longmapsto
\begin{cases}
\varphi_i(x \wedge \rho)_i = 0\\
\varphi_i(x \wedge \rho)_{\rho i} = x\\
\varphi_i(x \wedge \rho)_{\rho^2 i} = -x
\end{cases}
\]
\note{avant l'accolade, « $(x, \rho)$ » est biffé ; dans la première ligne, un indice surchargé.}

\marginal{\textbf{NB} \struck{$\varphi_i(x) + \varphi_{\rho i}(x) + \varphi_{\rho^2 i}(x) = 0$} \quad $\sum_{i \in I} \varphi_i(x) = 0$ \quad $\forall x \in V^\omega$. En fait, \ill{} sont exactes :
\[
0 \to V^\omega \xrightarrow{\ \mathrm{diag}\ } (V^\omega)^I \to \Gamma \to 0,
\]
i.e.
\[
0 \to V^\omega \to (V^\omega)^I \to V^I \to V \to 0.
\]}

On cherche à obtenir les systèmes $\psi$

\begin{tikzcd}
V^\omega \arrow[r, "\varphi_i"] \arrow[dr, "\psi_i"'] & \Gamma\\
 & \widetilde{\Gamma} \arrow[u, "u"']
\end{tikzcd}

avec

a) $\widetilde{\Gamma} \xrightarrow{u} \Gamma$ est un hom.\ de groupes : \uncertain{surjectif} \ill{} \uncertain{de noyau} $Z$ ;

b) $\psi_i$ un hom.\ $V^\omega \to \widetilde{\Gamma}$ qui relève $\varphi_i$ :
\[
u \psi_i = \varphi_i .
\]

\textbf{NB} $u$ \uncertain{donc} \ill{}, donc $\widetilde{\Gamma}$ est une \uncertain{extension} de $\Gamma$ par $Z$ (écrit additivement).

Soit $\rho \in \omega$, $i_0 \in I$, d'où $i_1 = \rho i_0$, $i_2 = \rho^2 i_0 = \rho i_1$. On aura
\[
\varphi_{i_0}(x \wedge \rho) + \varphi_{i_1}(x \wedge \rho) + \varphi_{i_2}(x \wedge \rho) = 0
\]
d'où
\[
\psi_{i_0}(x \wedge \rho)\, \psi_{i_1}(x \wedge \rho)\, \psi_{i_2}(x \wedge \rho) = z(x, \rho, i_0) \in Z .
\]
\note{l'argument se poursuit p.~33.}

\page{33}
On a, comme $Z$ est \uncertain{central},
\[
z(x, \rho, i_0) = z(x, \rho, i_1) = z(x, \rho, i_2),
\]
soit
\[
z = z(x, \rho) = z_\rho(x).
\]
D'ailleurs on aura
\[
\psi_{i_2}(x \wedge \rho)^{-1}\, \psi_{i_1}(x \wedge \rho)^{-1}\, \psi_{i_0}(x \wedge \rho)^{-1} = -z(x, \rho)
\]
i.e.
\[
\psi_{i_2}((-x) \wedge \rho)\, \psi_{i_1}((-x) \wedge \rho)\, \psi_{i_0}((-x) \wedge \rho) = -z(x, \rho)
\]
i.e. \struck{\uncertain{compte tenu}} comme $(-x) \wedge \rho = x \wedge \bar\rho$ \quad ($\bar\rho = \rho^2 = \rho^{-1}$)
\[
\psi_{i_2}(x \wedge \bar\rho)\, \psi_{i_1}(x \wedge \bar\rho)\, \psi_{i_0}(x \wedge \bar\rho) = -z(x, \rho)
\]
i.e.\ on a
\[
z(x, \bar\rho) = -z(x, \rho) \qquad \text{i.e.} \quad z_{\bar\rho} = -z_\rho
\]
\note{sous $z(x, \bar\rho)$, un petit signe, peut-être « $\shortparallel$ ».}

i.e.\ on a deux applications
\[
z_\rho : V \longrightarrow Z, \qquad \rho \in \omega
\]
\note{l'exposant de $V$ est surchargé (peut-être $V^\omega$ corrigé en $V$).}
\uncertain{échangées} par le signe (i.e.\ un élément de $(Z^V) \wedge_{\pm 1} \omega$). \struck{On aura}

\textbf{Prop.} On aura
\[
\begin{cases}
\text{\struck{$z_\rho(0) = 0$ et}} \quad \text{si } j = \rho i\\
\varphi_i(x)\, \varphi_j(y)\, \varphi_i(x)^{-1}\, \varphi_j(y)^{-1} = z_\rho(x) + z_\rho(y) - z_\rho(x + y)
\end{cases}
\]
\note{les $\varphi$ de cette formule sont peut-être des $\psi$ ; le tracé ne les distingue pas nettement.}

L'application $\lambda = z_\rho : V \to Z$ satisfait
\[
(*)\quad
\begin{cases}
z_\rho(0) = 0\\
z_\rho(x) + z_\rho(y + u) + z_\rho(y + v) + z_\rho(x + u + v)\\
\qquad = z_\rho(y) + z_\rho(x + u) + z_\rho(x + v) + z_\rho(y + u + v)
\end{cases}
\]
\note{la seconde équation est très surchargée : les lettres lues $u$, $v$ ont été récrites, et le premier terme du second membre pourrait aussi être autre chose ; lecture d'ensemble incertaine.}

\page{35}
\note{page écrite au crayon, d'un trait pâle ; elle reprend la relation $(*)$ de la p.~33. Le milieu de la page est biffé de plusieurs traits.}

ou encore (pour la $2^{\text{ème}}$ relation)
\[
\begin{aligned}
&\lambda(x + u + v) - \lambda(x + u) - \lambda(x + v) + \lambda(x)\\
&\qquad = \lambda(y + u + v) - \lambda(y + u) - \lambda(y + v) + \lambda(y)
\end{aligned}
\]
\[
(\Delta_{u,v} \lambda)(x) = \text{c}^{\text{te}} \quad \text{(pour $u, v$ fixés)}
\]
i.e. \struck{\ill{} \ill{}}, \uncertain{posant}
\[
\lambda(u + v) - \lambda(u) - \lambda(v) = \mu(u, v)
\]
on doit avoir
\[
\boxed{\lambda(x + u + v) - \lambda(x + u) - \lambda(x + v) + \lambda(x) = \mu(u, v)}
\]
i.e.\ $\mu$ est \emph{bi-additive}. \struck{$\lambda(x + u + v) = \lambda(u) + \lambda(v)$ (bilinéaire symétrique)}

Pour $\mu$ fixé, les $\lambda$ satisfaisant \add{et telles que $\lambda(0) = 0$} la relation précédente forment un torseur sous l'ens.\ des $f$ telles que
\[
\begin{cases}
f(0) = 0\\
f(x + u + v) - f(x + u) - f(x + v) + f(x) = 0
\end{cases}
\]
\note{le dernier terme, d'abord autre chose, est corrigé en $f(x)$.}

\struck{telles que \ldots\ $f(x) + f(y)$ \ldots\ i.e.\ $\exists\, g : V \to \ldots$ telle que $f(x+y) + f(0) = f(x) + f(y)$} \note{passage biffé de plusieurs traits ; lecture très incertaine.}

i.e.\ $f$ est un \struck{\uncertain{affine}} \emph{hom.}

Lorsque $V$ est un \uncertain{module} libre \struck{\ill{} si \ill{} qui} sur un $\mathbb{Z}/n\mathbb{Z}$, \uncertain{dans} \emph{le} \emph{\uncertain{produit}} \ill{} \uncertain{symétrique} de $V$ dans un $Z$ est \uncertain{associée} : une forme quadratique (\uncertain{définie} mod \uncertain{l'application} \emph{additive}) ssi les $\mu(e_i, e_i) \in 2Z$ ($(e_i)$ est une base de $V$ sur $\mathbb{Z}/n\mathbb{Z}$) \ldots
\note{la page s'arrête sur ces mots ; la phrase ne se poursuit pas sur la page suivante.}

\page{37}
Pour une application
\[
z_\rho : V \to Z
\]
de $V$ dans un groupe commutatif $Z$, satisfaisant les deux relations précédentes, $\exists$ une extension $\widetilde{\Gamma}^{(z_\rho)}$ de $\Gamma$ par $Z$, et des $\psi_i^{(z_\rho)}$ relevant les $\varphi_i$, donnant naissance (via $\rho$) à $z_\rho$, et le système de cette extension et des $\psi_i$ ($i \in I$) est déterminé à \uncertain{isom.}\ unique près.

Il y a un $Z$ \og universel\fg{} $\mathfrak{Z}_V$, savoir le quotient de $\mathbb{Z}^{(V)}$ par les \og relations\fg{} précédentes.\add{\marginal{écrit verticalement dans la marge gauche, et renvoyé ici par un signe : Il y a donc une \ill{} \uncertain{extension} universelle $\widetilde{\Gamma}$ de $\Gamma$ par un \ill{}, qui une fois \uncertain{choisi} $\rho$, \uncertain{le noyau} \ill{} \uncertain{indépendamment} des deux choix de $\rho$ \uncertain{possibles}, \ill{} $z_\rho : V \to \mathfrak{Z}^\omega$ \ill{}, est alors $\mathfrak{Z}^\omega = \mathfrak{Z} \wedge_{\pm 1} \omega$, i.e.\ $\lambda : V \to \mathfrak{Z}$ \ill{}, $z_\rho(x) = \lambda(x \wedge \rho)$ \ldots}} Notons que si $z_\rho$ est \struck{\ill{}} un hom.\ de groupes \uncertain{resp.} une application $2$-quadratique, les \uncertain{deux} relations \uncertain{sont} satisfaites, ce qui montre que $\mathfrak{Z}_V$ contient comme \ldots
\note{la phrase se poursuit en tête de la p.~38.}

\marginal{En haut de la marge gauche, écrit verticalement : et on aura bien
\[
z_{\bar\rho}(x) = \lambda(x \wedge \bar\rho) = \lambda(-x \wedge \rho) = -\lambda(x \wedge \rho) = -z_\rho(x),
\]
i.e.\ $z_{\bar\rho}(x) = -z_\rho(x)$ \ldots}

\page{38}
quotient\add{t} le $\mathbb{Z}$-module $V + \Gamma^2_{\mathbb{Z}}(V)$ (qui est le $\mathbb{Z}$-module universel qui reçoit $V$ par une application pol.\ de degré $\leq 2$ sans terme constant).
\marginal{dans la marge, renvoyé par un signe avant « le » : \uncertain{les} $z_\rho$ \uncertain{pris} dans \ldots\ \uncertain{sous-groupe} engendré par \uncertain{les} $x + \gamma^{(2)}$ \ldots\ $(x \in V)$.}

\textbf{Ex.} $V$ est un vectoriel \struck{de dim.~$2$} sur $\mathbb{F}_2$.
\struck{Soit $D = V^* = V \setminus \{0\}$, les droites de $V$ \ldots\ quelconques : celles de l'ens.\ $D$ \ldots\ et $\mathfrak{Z} = \mathbb{Z}^{(D)} / \text{relations}$.}
\note{passage barré de traits obliques et encadré, que la suite remplace.}

Dans ce cas, \struck{\uncertain{puisque}} \add{\uncertain{vu}} que les applications $\lambda : V \to Z$ telles que
\[
\lambda(x + y) - \lambda(x) - \lambda(y) = \mu(x, y) \qquad \text{($\mu$ $\mathbb{Z}$-bil.\ sym.)}
\]
sont polynomiales (en effet, on aura $\lambda(0) = 0 \Rightarrow$ \uncertain{d'où}
\[
\lambda(2x) = 2\lambda(x) + \mu(x, x),
\]
\uncertain{donc} $\mu(x, x) = -2\lambda(x)$, \uncertain{donc} $\mu(x, x) \in 2Z$ $(\forall x \in V)$, d'où \uncertain{facilement} le résultat). Donc
\[
\mathfrak{Z} = V + \Gamma^2_{\mathbb{Z}}(V) \simeq (\mathbb{Z}/4\mathbb{Z})^{I} + (\mathbb{Z}/2\mathbb{Z})^{P_2(I)}
\]
\note{sous $\lambda(2x)$, un « $\shortparallel 0$ » ($2x = 0$ sur $\mathbb{F}_2$). Dans la dernière ligne, après $(\mathbb{Z}/4\mathbb{Z})^I$, un « $+\,\mathbb{Z}/4\mathbb{Z}$ » est biffé ; l'accolade sous $\Gamma^2_{\mathbb{Z}}(V)$ porte « $\shortparallel$ ».}

\drawing{dans la marge gauche, un petit schéma : deux paires de points $x, y$ et $x', y'$ reliées par des flèches courbes en sens opposés, autour d'un axe vertical.}

T.S.V.P.

\page{39}
\note{suite de la p.~38 (« T.S.V.P. »).}

Supposons que $V$ soit de dimension $2$ \add{sur $\mathbb{F}_2$}. Alors il y a une forme quadratique canonique sur $V$, définie par
\[
q(x) =
\begin{cases}
0 & \text{si } x = 0\\
1 & \text{si } x \neq 0
\end{cases}
\]
qui donne donc lieu à une ext.\ centrale canonique de $\Gamma(V)$ par $\mathbb{Z}/2\mathbb{Z} = \mathbb{F}_2$ telles que
\[
\psi_{i_0}(x)\, \psi_{i_1}(x)\, \psi_{i_2}(x) = z \qquad \text{(él.\ non nul de $Z = \mathbb{Z}/2\mathbb{Z}$)}
\]
$\forall x \in V^\omega$ (et un \struck{\uncertain{ordre} \uncertain{circulaire}} \add{\uncertain{triplet}} $(i_0, i_1, i_2)$ \struck{\ill{}} sur $I$ \ldots)
\note{la page s'arrête ici ; le reste est blanc.}

\end{document}
